\chapter{FCC Contact Microphysics and Long-Wavelength Homogenization}
\label{app:fcc_microphysics}

\section{Scope, Status of Results, and Constitutive Limits}

This appendix develops the leading reduced contact mechanics of the active FCC Soft Floor. Its purpose is to connect the local endpoint-transfer law used in the gravitational propagation theory to an explicit conservative contact action, objective nearest-neighbor kinematics, and the long-wavelength acoustic tensor of the FCC lattice.

The analysis is restricted to the active FCC phase,
\begin{equation}
    0 \leq \alpha_{\mathrm{G}} < 1,
\end{equation}
where $\alpha_{\mathrm{G}}$ denotes the dimensionless gravitational loading fraction of the local active capacity. The limiting value ($\alpha_{\mathrm{G}}=1$) is not treated as an ordinary elastic state. It marks termination of the active FCC coordinates and transition toward the phase-locked SC substrate state.

The gravitational background is assumed quasistatic over one local transfer event,
\begin{equation}
    \dot{\alpha}_{\mathrm{G}} \approx 0,
\end{equation}
so that the endpoint-transfer factor acts as a time-independent holonomic constraint during the local oscillation. The coherent transverse sector is considered first. The interstitial-condensate pressure constraint and localized defect-induced relative-flow sector are separated into the following appendix.

The present appendix does not claim to derive the endpoint-transfer product from a complete parent node--condensate--Tension-Column potential. Instead, it gives a reduced work-conjugate construction of that endpoint law and proves that, once imposed as the local transfer constraint, the law is compatible with a conservative action, a positive contact Hessian, and the long-wavelength transverse propagation rule
\begin{equation}
    c_P = c_{\mathrm{E}}(1-\alpha_{\mathrm{G}})^2.
\end{equation}

This distinction is essential. The endpoint product is treated as a reduced constitutive closure. The constrained kinetic action then determines its inertial and dispersive consequences.

The mechanical hierarchy developed here is
\begin{equation}
\begin{array}{c}
    \text{FCC contact geometry} \\
    \longrightarrow \text{objective bond kinematics} \\
    \longrightarrow \text{endpoint-transfer closure} \\
    \longrightarrow \text{referred inertia} \\
    \longrightarrow \text{long-wavelength acoustic tensor}.
\end{array}
\end{equation}

The full finite-wave-number Bloch spectrum of the active FCC lattice is not completed in this appendix. A formal stiffness assembly is given, but the explicit full dynamic matrix, its eigenvalue branches over the first Brillouin zone, and the complete stability proof remain open microscopic calculations.

The status of the main results is as follows.

\begin{table}[h!]
\centering
\renewcommand{\arraystretch}{1.3}
\begin{tabular}{p{0.45\linewidth} p{0.5\linewidth}}
\hline
\textbf{Result} & \textbf{Status in this appendix} \\
\hline
$d_{\mathrm{nn}}=2r_{\mathrm{D}}$ & FCC geometric identity \\
$V_{\mathrm{p}}=4\sqrt{2}r_{\mathrm{D}}^3$ & primitive-cell geometric identity \\
endpoint product $\mathcal{C}_{ij}=(1-\alpha_i)(1-\alpha_j)$ & reduced work-conjugate constitutive closure \\
$\mu_{\mathrm{eff}}=\mu_{\mathrm{D}}/\mathcal{C}^2$ & consequence of constrained kinetic reduction \\
$c_P=c_{\mathrm{E}}\mathcal{C}$ & long-wavelength consequence of the reduced action \\
objective contact measures $\delta_{\mathrm{n}}, \boldsymbol{\delta}_{\mathrm{t}}, \boldsymbol{\delta}_{\mathrm{b}}, \delta_{\mathrm{tw}}, \Delta\varphi$ & derived from frame-indifferent contact kinematics \\
parity-even quadratic contact action & leading local bond model \\
FCC cubic elastic constants & derived in the affine ($k\rightarrow0$) limit \\
$\kappa_{\mathrm{n}}=\kappa_{\mathrm{t}}$ & condition for exact long-wavelength isotropy \\
full Bloch dynamic matrix & defined formally, not fully calculated \\
microscopic-to-macroscopic modulus bridge & identified as a required coarse-graining problem \\
\hline
\end{tabular}
\end{table}

The active Soft Floor is treated as a dual mechanical structure. The solid FCC scaffold supplies local contact rigidity. The interstitial condensate supplies entrained inertial participation and, in the coherent branch, supports the constraint structure developed separately in Appendix \ref{app:condensate_impedance}. The Hard Floor does not appear here as an ordinary finite-frequency oscillator. It enters as the singular vertical constraint that limits the admissible active transfer through $\alpha_{\mathrm{G}}$.

The contact coefficients introduced below are tangent coefficients about an active-phase equilibrium. They are not assumed to remain valid through the nonlinear saturation transition at $\alpha_{\mathrm{G}}=1$. Throughout this appendix, all small-strain contact statements are therefore understood in the active-side limit
\begin{equation}
    \alpha_{\mathrm{G}} \rightarrow 1^-,
\end{equation}
not at or beyond the phase boundary.

\section{Constrained Transverse FCC Contact Action}

Consider a representative nearest-neighbor transfer path in the active FCC Soft Floor. The nearest-neighbor separation is
\begin{equation}
    d_{\mathrm{nn}} = 2r_{\mathrm{D}}.
\end{equation}

Let $u_n(t)$ denote the completed transverse displacement of the $n$-th active FCC node along one selected polarization orthogonal to the bond direction. For the contact between neighboring nodes $n$ and $n+1$, the transverse strain is
\begin{equation}
    \varepsilon_{n,n+1}^{\perp} = \frac{u_{n+1}-u_n}{d_{\mathrm{nn}}}.
\end{equation}

The baseline active surface shear tension is $\tau_{\mathrm{D}}$. The transverse contact energy per unit participating interface area is therefore
\begin{equation}
    \mathcal{U}_{n,n+1}^{\perp} = \frac{\tau_{\mathrm{D}}}{2d_{\mathrm{nn}}^2} \left(u_{n+1}-u_n\right)^2.
\end{equation}

The corresponding two-node potential Hessian is
\begin{equation}
    \mathbf{H}_U^{(n,n+1)} = \frac{\tau_{\mathrm{D}}}{d_{\mathrm{nn}}^2}
    \begin{pmatrix}
        1 & -1 \\
        -1 & 1
    \end{pmatrix}.
\end{equation}

At the retained active-phase order, gravitational loading does not continuously soften this transverse potential Hessian. Instead,
\begin{equation}
    \mathbf{H}_U^{(n,n+1)}(\alpha_{\mathrm{G}}) = \mathbf{H}_U^{(n,n+1)}(0), \qquad 0 \leq \alpha_{\mathrm{G}} < 1.
\end{equation}

The gravitational dependence enters through the kinetic sector by reducing the completed transfer relative to the underlying condensate-supported reconfiguration.

Let $y_n(t)$ denote the latent local reconfiguration coordinate required to produce the completed displacement $u_n(t)$. In the unstrained active vacuum,
\begin{equation}
    \alpha_{\mathrm{G}} = 0, \qquad u_n = y_n.
\end{equation}

For a locally uniform gravitational strain, the endpoint-transfer factor is
\begin{equation}
    \mathcal{C} = (1-\alpha_{\mathrm{G}})^2.
\end{equation}

The completed displacement is constrained to the latent reconfiguration by
\begin{equation}
    u_n = \mathcal{C} y_n.
\end{equation}

Equivalently,
\begin{equation}
    y_n = \frac{u_n}{\mathcal{C}}.
\end{equation}

The factor $\mathcal{C}$ is therefore the local transfer Jacobian between the underlying reconfiguration coordinate and the completed active-FCC displacement.

The baseline surface inertia is $\mu_{\mathrm{D}}$. The kinetic energy per unit participating interface area associated with the latent coordinate is
\begin{equation}
    \mathcal{T}_n = \frac{1}{2}\mu_{\mathrm{D}}\dot{y}_n^2.
\end{equation}

A reduced directional Lagrangian per unit participating interface area is
\begin{equation}
    \mathscr{L}_\perp = \sum_n \left[ \frac{1}{2}\mu_{\mathrm{D}}\dot{y}_n^2 - \frac{\tau_{\mathrm{D}}}{2d_{\mathrm{nn}}^2} \left(u_{n+1}-u_n\right)^2 \right],
\end{equation}
subject to the holonomic constraint
\begin{equation}
    u_n - \mathcal{C} y_n = 0.
\end{equation}

Since the background is quasistatic,
\begin{equation}
    \dot{\mathcal{C}} = 0,
\end{equation}
and the constraint gives
\begin{equation}
    \dot{y}_n = \frac{\dot{u}_n}{\mathcal{C}}.
\end{equation}

Eliminating $y_n$ gives the reduced transverse Lagrangian
\begin{equation}
    \mathscr{L}_\perp^{(\mathrm{red})} = \sum_n \left[ \frac{1}{2} \frac{\mu_{\mathrm{D}}}{\mathcal{C}^2} \dot{u}_n^2 - \frac{\tau_{\mathrm{D}}}{2d_{\mathrm{nn}}^2} \left(u_{n+1}-u_n\right)^2 \right].
\end{equation}

Thus the effective surface inertia in the completed coordinate is
\begin{equation}
    \mu_{\mathrm{eff}} = \frac{\mu_{\mathrm{D}}}{\mathcal{C}^2}.
\end{equation}

Using $\mathcal{C} = (1-\alpha_{\mathrm{G}})^2$, one obtains
\begin{equation}
    \boxed{ \mu_{\mathrm{eff}} = \frac{\mu_{\mathrm{D}}}{(1-\alpha_{\mathrm{G}})^4} }.
\end{equation}

The inverse-fourth dependence is not imposed independently. It is the kinetic consequence of two operations: a two-endpoint transfer reduction and a quadratic kinetic-energy transformation.

Varying the reduced action with respect to $u_n$ gives
\begin{equation}
    \frac{\mu_{\mathrm{D}}}{\mathcal{C}^2}\ddot{u}_n = \frac{\tau_{\mathrm{D}}}{d_{\mathrm{nn}}^2} \left(u_{n+1}-2u_n+u_{n-1}\right).
\end{equation}

Equivalently,
\begin{equation}
    \mu_{\mathrm{eff}}\ddot{u}_n = \frac{\tau_{\mathrm{D}}}{d_{\mathrm{nn}}^2} \left(u_{n+1}-2u_n+u_{n-1}\right).
\end{equation}

For a plane-wave mode,
\begin{equation}
    u_n(t) = u_0 \exp\left[i(knd_{\mathrm{nn}}-\omega t)\right],
\end{equation}
the nearest-neighbor difference becomes
\begin{equation}
    u_{n+1}-2u_n+u_{n-1} = -4u_n \sin^2\left(\frac{kd_{\mathrm{nn}}}{2}\right).
\end{equation}

Substitution gives the directional dispersion relation
\begin{equation}
    \omega^2(k,\alpha_{\mathrm{G}}) = \frac{4\tau_{\mathrm{D}}}{\mu_{\mathrm{D}}d_{\mathrm{nn}}^2} \mathcal{C}^2 \sin^2\left(\frac{kd_{\mathrm{nn}}}{2}\right).
\end{equation}

The baseline active-vacuum identity is
\begin{equation}
    c_{\mathrm{E}}^2 = \frac{\tau_{\mathrm{D}}}{\mu_{\mathrm{D}}}.
\end{equation}

Therefore,
\begin{equation}
    \boxed{ \omega^2(k,\alpha_{\mathrm{G}}) = \frac{4c_{\mathrm{E}}^2\mathcal{C}^2}{d_{\mathrm{nn}}^2} \sin^2\left(\frac{kd_{\mathrm{nn}}}{2}\right) }.
\end{equation}

In the local continuum limit, $kd_{\mathrm{nn}} \ll 1$, so
\begin{equation}
    \sin\left(\frac{kd_{\mathrm{nn}}}{2}\right) = \frac{kd_{\mathrm{nn}}}{2} + \mathcal{O}(k^3d_{\mathrm{nn}}^3).
\end{equation}

The dispersion reduces to
\begin{equation}
    \omega^2 = c_{\mathrm{E}}^2\mathcal{C}^2 k^2 + \mathcal{O}(k^4d_{\mathrm{nn}}^2).
\end{equation}

Hence,
\begin{equation}
    \lim_{k\rightarrow0}\frac{\omega}{k} = c_{\mathrm{E}}\mathcal{C}.
\end{equation}

Using the local endpoint factor, the long-wavelength phase velocity is
\begin{equation}
    \boxed{ c_P = c_{\mathrm{E}}(1-\alpha_{\mathrm{G}})^2 }.
\end{equation}

The positive-frequency branch may be written as
\begin{equation}
    \omega(k,\alpha_{\mathrm{G}}) = \frac{2c_{\mathrm{E}}\mathcal{C}}{d_{\mathrm{nn}}} \left| \sin\left(\frac{kd_{\mathrm{nn}}}{2}\right) \right|.
\end{equation}

Within the first monotonic branch, the group velocity is
\begin{equation}
    v_g = \frac{\mathrm{d}\omega}{\mathrm{d}k} = c_{\mathrm{E}}\mathcal{C} \cos\left(\frac{kd_{\mathrm{nn}}}{2}\right).
\end{equation}

Therefore,
\begin{equation}
    \lim_{k\rightarrow0}v_g = c_{\mathrm{E}}\mathcal{C} = c_P.
\end{equation}

The same constraint rescales both the long-wavelength phase velocity and the long-wavelength group velocity.

As the active-phase boundary is approached ($\alpha_{\mathrm{G}} \rightarrow 1^-$), the transfer factor vanishes ($\mathcal{C} \rightarrow 0$), and the referred inertia diverges:
\begin{equation}
    \mu_{\mathrm{eff}} = \frac{\mu_{\mathrm{D}}}{\mathcal{C}^2} \rightarrow \infty.
\end{equation}

For fixed lattice wave number, $\omega(k,\alpha_{\mathrm{G}}) \rightarrow 0$, and the continuum propagation capacity satisfies $c_P \rightarrow 0$. 

This limiting behavior does not mean that the SC substrate has infinite material mass. It means that the completed active-FCC coordinate loses admissibility at the phase boundary. The FCC mode is arrested by divergent referred inertia before the active coordinates terminate.

This one-dimensional contact sequence is only a directional reduction. It is used to expose the inertial consequence of the endpoint-transfer constraint. In the full FCC assembly, the six undirected nearest-neighbor bonds contribute stiffness to the primitive-cell potential, while translational, microrotational, phase, and condensate inertia are assigned once through the primitive-cell mass matrix. The directional derivation must therefore not be interpreted as assigning an independent complete nodal mass to every bond.

\section{Work-Conjugate Construction of the Endpoint Transfer Factor}

The preceding section used the local transfer coefficient $\mathcal{C}$ to express the constrained relation between the latent node--condensate reconfiguration and the completed active-FCC displacement. This section gives the reduced work-conjugate construction of that coefficient. The construction is not presented as a complete derivation from a parent node--condensate--Tension-Column potential. It defines the reduced endpoint closure that such a parent theory must reproduce.

Consider one endpoint $i$ of a nearest-neighbor contact. The local gravitational loading fraction is
\begin{equation}
    \alpha_{\mathrm{G}}^{(i)} = \frac{\tau_L^{(i)}}{\tau_{\mathrm{D}}},
\end{equation}
where $\tau_L^{(i)}$ is the local mass-induced loading of the active Soft Floor and $\tau_{\mathrm{D}}$ is the baseline active surface-tension scale. The residual active capacity available for additional local transfer is
\begin{equation}
    \tau_{\mathrm{av}}^{(i)} = \tau_{\mathrm{D}} - \tau_L^{(i)}.
\end{equation}

Therefore,
\begin{equation}
    \tau_{\mathrm{av}}^{(i)} = \tau_{\mathrm{D}} \left( 1-\alpha_{\mathrm{G}}^{(i)} \right).
\end{equation}

The dimensionless single-endpoint transfer factor is defined by
\begin{equation}
    \mathcal{C}_i \equiv 1-\alpha_{\mathrm{G}}^{(i)}.
\end{equation}

Thus,
\begin{equation}
    \tau_{\mathrm{av}}^{(i)} = \tau_{\mathrm{D}}\mathcal{C}_i.
\end{equation}

Inside the active phase,
\begin{equation}
    0 < \mathcal{C}_i \leq 1.
\end{equation}

At zero gravitational loading,
\begin{equation}
    \alpha_{\mathrm{G}}^{(i)}=0, \qquad \mathcal{C}_i=1.
\end{equation}

At the active-side saturation boundary,
\begin{equation}
    \alpha_{\mathrm{G}}^{(i)}\rightarrow1^{-}, \qquad \mathcal{C}_i\rightarrow0.
\end{equation}

Let $y_i$ denote the latent endpoint reconfiguration coordinate, and let $u_i$ denote the completed displacement accepted by the active FCC contact state. Introduce the dimensionless coordinates
\begin{equation}
    \eta_i = \frac{y_i}{d_{\mathrm{nn}}}, \qquad \varepsilon_i = \frac{u_i}{d_{\mathrm{nn}}}.
\end{equation}

The latent coordinate $\eta_i$ measures the node--condensate reconfiguration presented to the endpoint. The completed coordinate $\varepsilon_i$ measures the displacement transmitted into the active FCC contact state.

The reduced endpoint closure is the equality of the incremental work represented in the latent and completed coordinate descriptions:
\begin{equation}
    \tau_{\mathrm{D}} \delta\varepsilon_i = \tau_{\mathrm{av}}^{(i)} \delta\eta_i.
\end{equation}

Substituting the residual capacity gives
\begin{equation}
    \tau_{\mathrm{D}} \delta\varepsilon_i = \tau_{\mathrm{D}}\mathcal{C}_i \delta\eta_i.
\end{equation}

Hence,
\begin{equation}
    \delta\varepsilon_i = \mathcal{C}_i \delta\eta_i.
\end{equation}

For a quasistatic gravitational background, $\mathcal{C}_i$ is constant during the local transfer. Integration from the common undeformed reference state gives
\begin{equation}
    \varepsilon_i = \mathcal{C}_i\eta_i.
\end{equation}

Multiplying by $d_{\mathrm{nn}}$, one obtains
\begin{equation}
    u_i = \mathcal{C}_i y_i.
\end{equation}

Thus $\mathcal{C}_i$ is the single-endpoint transfer factor between the latent reconfiguration coordinate and the completed active-FCC displacement. It is not the Scale-Invariant Mass Moment $J_i$. The latter remains reserved throughout TFT for
\begin{equation}
    J_i = \rho_i r_i^4.
\end{equation}

For a nearest-neighbor transfer between endpoints $i$ and $j$, the source endpoint releases only the fraction permitted by its residual active capacity:
\begin{equation}
    y_{ij}^{(\mathrm{rel})} = \mathcal{C}_i y_{ij}.
\end{equation}

The target endpoint then accepts the released amplitude according to its own residual capacity:
\begin{equation}
    u_{ij} = \mathcal{C}_j y_{ij}^{(\mathrm{rel})}.
\end{equation}

Therefore,
\begin{equation}
    u_{ij} = \mathcal{C}_i\mathcal{C}_j y_{ij}.
\end{equation}

The two-endpoint transfer factor is
\begin{equation}
    \mathcal{C}_{ij} \equiv \mathcal{C}_i\mathcal{C}_j.
\end{equation}

Using the endpoint definitions,
\begin{equation}
    \boxed{ \mathcal{C}_{ij} = \left( 1-\alpha_{\mathrm{G}}^{(i)} \right) \left( 1-\alpha_{\mathrm{G}}^{(j)} \right) }.
\end{equation}

The product form expresses sequential endpoint acceptance. The target endpoint acts on the amplitude released by the source endpoint, not on the original latent amplitude. The transfer sequence is
\begin{equation}
    y_{ij} \longrightarrow \mathcal{C}_i y_{ij} \longrightarrow \mathcal{C}_j\mathcal{C}_i y_{ij}.
\end{equation}

Because the endpoint factors are scalars in the selected transverse transfer sector,
\begin{equation}
    \mathcal{C}_{ij} = \mathcal{C}_{ji}.
\end{equation}

This reciprocity is required for a conservative nearest-neighbor contact in a static, unpolarized active vacuum.

For a gravitational field varying negligibly across one nearest-neighbor separation,
\begin{equation}
    \alpha_{\mathrm{G}}^{(i)} \approx \alpha_{\mathrm{G}}^{(j)} \approx \alpha_{\mathrm{G}}.
\end{equation}

Then
\begin{equation}
    \mathcal{C}_i \approx \mathcal{C}_j \approx 1-\alpha_{\mathrm{G}},
\end{equation}
and
\begin{equation}
    \mathcal{C}_{ij} \approx (1-\alpha_{\mathrm{G}})^2.
\end{equation}

The local continuum endpoint factor is therefore
\begin{equation}
    \mathcal{C} = (1-\alpha_{\mathrm{G}})^2.
\end{equation}

The approximation has a controlled endpoint-gradient structure. Define
\begin{equation}
    \alpha_{\mathrm{G}} = \frac{ \alpha_{\mathrm{G}}^{(i)} + \alpha_{\mathrm{G}}^{(j)} }{2},
\end{equation}
and
\begin{equation}
    \delta\alpha_{\mathrm{G}} = \alpha_{\mathrm{G}}^{(j)} - \alpha_{\mathrm{G}}^{(i)}.
\end{equation}

Then
\begin{equation}
    \mathcal{C}_{ij} = \left( 1-\alpha_{\mathrm{G}}^{(i)} \right) \left( 1-\alpha_{\mathrm{G}}^{(j)} \right)
\end{equation}
becomes
\begin{equation}
    \mathcal{C}_{ij} = (1-\alpha_{\mathrm{G}})^2 - \frac{ (\delta\alpha_{\mathrm{G}})^2 }{4}.
\end{equation}

The first-order endpoint-gradient correction cancels. The leading correction to the local continuum factor is second order in the loading difference across the bond.

The same transfer factor determines the referred inertia of the completed coordinate. For the latent coordinate,
\begin{equation}
    T = \frac{1}{2} \mu_{\mathrm{D}}\dot{y}^2.
\end{equation}

Using
\begin{equation}
    u = \mathcal{C} y,
\end{equation}
and assuming a quasistatic background,
\begin{equation}
    \dot{\mathcal{C}} = 0,
\end{equation}
one obtains
\begin{equation}
    \dot{y} = \frac{\dot{u}}{\mathcal{C}}.
\end{equation}

Therefore,
\begin{equation}
    T = \frac{1}{2} \frac{\mu_{\mathrm{D}}}{\mathcal{C}^2} \dot{u}^2.
\end{equation}

The effective surface inertia of the completed coordinate is
\begin{equation}
    \mu_{\mathrm{eff}} = \frac{\mu_{\mathrm{D}}}{\mathcal{C}^2}.
\end{equation}

For the local continuum endpoint factor, $\mathcal{C}=(1-\alpha_{\mathrm{G}})^2$, this becomes
\begin{equation}
    \boxed{ \mu_{\mathrm{eff}} = \frac{\mu_{\mathrm{D}}}{(1-\alpha_{\mathrm{G}})^4} }.
\end{equation}

The baseline transverse normalization is
\begin{equation}
    c_{\mathrm{E}}^2 = \frac{\tau_{\mathrm{D}}}{\mu_{\mathrm{D}}}.
\end{equation}

Since the active-phase contact Hessian is unchanged at the retained tangent order, while the completed-coordinate inertia is referred by $\mathcal{C}^{-2}$, the transverse propagation speed satisfies
\begin{equation}
    c_P^2 = \frac{\tau_{\mathrm{D}}}{\mu_{\mathrm{eff}}} = \frac{\tau_{\mathrm{D}}}{\mu_{\mathrm{D}}} \mathcal{C}^2.
\end{equation}

Hence,
\begin{equation}
    c_P^2 = c_{\mathrm{E}}^2\mathcal{C}^2.
\end{equation}

Taking the positive propagation branch gives
\begin{equation}
    c_P = c_{\mathrm{E}}\mathcal{C}.
\end{equation}

Using $\mathcal{C}=(1-\alpha_{\mathrm{G}})^2$, one obtains
\begin{equation}
    \boxed{ c_P = c_{\mathrm{E}} (1-\alpha_{\mathrm{G}})^2 }.
\end{equation}

The reduced endpoint chain is therefore
\begin{equation}
    \alpha_{\mathrm{G}}^{(i)} \longrightarrow \mathcal{C}_i \longrightarrow \mathcal{C}_{ij} \longrightarrow \mu_{\mathrm{eff}} \longrightarrow c_P.
\end{equation}

This section establishes the reduced work-conjugate endpoint closure used by the active-FCC transverse action. A complete microscopic potential must still derive $\mathcal{C}_i$ from the coupled normal, transverse, phase, condensate, and vertical Tension-Column degrees of freedom. The present result specifies the conservative endpoint map and its inertial consequence within the active phase.

\section{Objective Nearest-Neighbor FCC Contact Kinematics}

The directional action captures the transverse transfer law along one selected nearest-neighbor sequence. To assemble the active Soft Floor as a three-dimensional FCC lattice, the contact law must be written in an objective form. A contact energy depending only on the relative displacement of two node centers is not sufficient, because a rigid-body rotation produces relative center motion even when the physical contact surfaces undergo no elastic deformation. The bond kinematics must therefore include nodal translations, nodal microrotations, and internal phase differences.

Consider two nearest-neighbor FCC nodes labeled $i$ and $j$. Their reference center-to-center vector is
\begin{equation}
    \mathbf{d}_{ij} = d_{\mathrm{nn}}\hat{\mathbf{n}}_{ij},
\end{equation}
where
\begin{equation}
    d_{\mathrm{nn}} = 2r_{\mathrm{D}}, \qquad |\hat{\mathbf{n}}_{ij}| = 1.
\end{equation}

The normal projector associated with the contact is
\begin{equation}
    \mathbf{P}_{\mathrm{n}}^{(ij)} = \hat{\mathbf{n}}_{ij}\hat{\mathbf{n}}_{ij}^{\mathsf{T}},
\end{equation}
and the transverse projector is
\begin{equation}
    \mathbf{P}_{\perp}^{(ij)} = \mathbf{I} - \hat{\mathbf{n}}_{ij}\hat{\mathbf{n}}_{ij}^{\mathsf{T}}.
\end{equation}

Let $\mathbf{u}_i, \mathbf{u}_j$ be the infinitesimal translations of the two node centers. Let $\boldsymbol{\vartheta}_i, \boldsymbol{\vartheta}_j$ be their infinitesimal microrotations. These are mechanical rotations of the local contact frames. They are not identical to the internal scalar phase variable. The internal phases are denoted $\varphi_i, \varphi_j$.

Thus the contact theory uses three types of local variables:
\begin{itemize}
    \item $\mathbf{u}$ for center translation,
    \item $\boldsymbol{\vartheta}$ for mechanical microrotation,
    \item $\varphi$ for internal topological phase.
\end{itemize}
They are kept distinct at the present constitutive order.

The reference contact point lies at $+r_{\mathrm{D}}\hat{\mathbf{n}}_{ij}$ relative to the center of node $i$, and at $-r_{\mathrm{D}}\hat{\mathbf{n}}_{ij}$ relative to the center of node $j$. The displacement of the contact point carried by node $i$ is
\begin{equation}
    \mathbf{u}_i^{\mathrm{c}} = \mathbf{u}_i + r_{\mathrm{D}} \boldsymbol{\vartheta}_i \times \hat{\mathbf{n}}_{ij}.
\end{equation}

The displacement of the corresponding contact point carried by node $j$ is
\begin{equation}
    \mathbf{u}_j^{\mathrm{c}} = \mathbf{u}_j - r_{\mathrm{D}} \boldsymbol{\vartheta}_j \times \hat{\mathbf{n}}_{ij}.
\end{equation}

The relative displacement of the two material contact points is therefore $\Delta\mathbf{u}_{ij}^{\mathrm{c}} = \mathbf{u}_j^{\mathrm{c}} - \mathbf{u}_i^{\mathrm{c}}$. Substitution gives
\begin{equation}
    \boxed{ \Delta\mathbf{u}_{ij}^{\mathrm{c}} = \mathbf{u}_j - \mathbf{u}_i - r_{\mathrm{D}} (\boldsymbol{\vartheta}_i+\boldsymbol{\vartheta}_j) \times \hat{\mathbf{n}}_{ij} }.
\end{equation}
This is the objective contact-point displacement. It is the correct local deformation measure entering the nearest-neighbor contact energy.

The normal contact deformation is
\begin{equation}
    \delta_{\mathrm{n}}^{(ij)} = \hat{\mathbf{n}}_{ij} \cdot \Delta\mathbf{u}_{ij}^{\mathrm{c}}.
\end{equation}
The rotational term is transverse to $\hat{\mathbf{n}}_{ij}$, and therefore
\begin{equation}
    \boxed{ \delta_{\mathrm{n}}^{(ij)} = \hat{\mathbf{n}}_{ij} \cdot (\mathbf{u}_j-\mathbf{u}_i) }.
\end{equation}

The objective tangential slip is
\begin{equation}
    \boldsymbol{\delta}_{\mathrm{t}}^{(ij)} = \mathbf{P}_{\perp}^{(ij)} \Delta\mathbf{u}_{ij}^{\mathrm{c}}.
\end{equation}
Explicitly,
\begin{equation}
    \boxed{ \boldsymbol{\delta}_{\mathrm{t}}^{(ij)} = \mathbf{P}_{\perp}^{(ij)} \left[ \mathbf{u}_j-\mathbf{u}_i - r_{\mathrm{D}} (\boldsymbol{\vartheta}_i+\boldsymbol{\vartheta}_j) \times \hat{\mathbf{n}}_{ij} \right] }.
\end{equation}
This is the relative displacement of the physical contact surfaces in the tangent plane. It differs from the simpler expression $\mathbf{P}_{\perp}^{(ij)}(\mathbf{u}_j-\mathbf{u}_i)$, which would incorrectly assign elastic slip to a rigidly rotating pair of nodes.

Define the relative microrotation by $\Delta\boldsymbol{\vartheta}_{ij} = \boldsymbol{\vartheta}_j-\boldsymbol{\vartheta}_i$. Its transverse part gives the local rolling or bending measure,
\begin{equation}
    \boxed{ \boldsymbol{\delta}_{\mathrm{b}}^{(ij)} = \mathbf{P}_{\perp}^{(ij)} \Delta\boldsymbol{\vartheta}_{ij} }.
\end{equation}
Its axial part gives the relative twist about the contact direction,
\begin{equation}
    \boxed{ \delta_{\mathrm{tw}}^{(ij)} = \hat{\mathbf{n}}_{ij} \cdot \Delta\boldsymbol{\vartheta}_{ij} }.
\end{equation}

The internal phase difference across the bond is
\begin{equation}
    \boxed{ \Delta\varphi_{ij} = \varphi_j-\varphi_i }.
\end{equation}
Only $\Delta\varphi_{ij}$, not $\varphi_i$ or $\varphi_j$ separately, may enter the contact energy. The bond action is therefore invariant under a uniform internal phase displacement,
\begin{equation}
    \varphi_i\rightarrow\varphi_i+\varphi_0, \qquad \varphi_j\rightarrow\varphi_j+\varphi_0.
\end{equation}
This is the local contact expression of the phase-shift freedom of the Soft Floor.

The objectivity of the construction can be checked directly. Consider a rigid-body motion composed of a uniform translation $\mathbf{a}$ and an infinitesimal rotation $\boldsymbol{\Omega}$. The nodal center displacements are
\begin{equation}
    \mathbf{u}_i = \mathbf{a}+\boldsymbol{\Omega}\times\mathbf{X}_i, \qquad \mathbf{u}_j = \mathbf{a}+\boldsymbol{\Omega}\times\mathbf{X}_j.
\end{equation}
Since $\mathbf{X}_j-\mathbf{X}_i = d_{\mathrm{nn}}\hat{\mathbf{n}}_{ij}$, the relative center displacement is
\begin{equation}
    \mathbf{u}_j-\mathbf{u}_i = d_{\mathrm{nn}} \boldsymbol{\Omega}\times\hat{\mathbf{n}}_{ij}.
\end{equation}

For a rigid rotation of the node frames, $\boldsymbol{\vartheta}_i = \boldsymbol{\vartheta}_j = \boldsymbol{\Omega}$. Then
\begin{equation}
    \Delta\mathbf{u}_{ij}^{\mathrm{c}} = d_{\mathrm{nn}} \boldsymbol{\Omega}\times\hat{\mathbf{n}}_{ij} - 2r_{\mathrm{D}} \boldsymbol{\Omega}\times\hat{\mathbf{n}}_{ij}.
\end{equation}
Using $d_{\mathrm{nn}}=2r_{\mathrm{D}}$, one obtains $\Delta\mathbf{u}_{ij}^{\mathrm{c}} = \mathbf{0}$.

Consequently,
\begin{equation}
    \delta_{\mathrm{n}}^{(ij)}=0, \qquad \boldsymbol{\delta}_{\mathrm{t}}^{(ij)}=\mathbf{0}.
\end{equation}
Because the two nodes undergo the same microrotation,
\begin{equation}
    \boldsymbol{\delta}_{\mathrm{b}}^{(ij)}=\mathbf{0}, \qquad \delta_{\mathrm{tw}}^{(ij)}=0.
\end{equation}
A uniform phase displacement gives $\Delta\varphi_{ij}=0$.

Thus all retained contact deformation measures vanish under rigid translation, rigid rotation, and uniform internal phase displacement. The bond law constructed from these measures is frame-indifferent to the linearized order retained here.

The resulting objective nearest-neighbor deformation set is
\begin{equation}
    \left\{ \delta_{\mathrm{n}}^{(ij)}, \boldsymbol{\delta}_{\mathrm{t}}^{(ij)}, \boldsymbol{\delta}_{\mathrm{b}}^{(ij)}, \delta_{\mathrm{tw}}^{(ij)}, \Delta\varphi_{ij} \right\}.
\end{equation}
This set is the local kinematic foundation for the parity-even quadratic bond action.

\section{Parity-Even Quadratic Bond Action}

The contact energy is now written as the most general baseline quadratic form built from the objective contact measures above, subject to three restrictions:
\begin{itemize}
    \item rigid-motion invariance,
    \item reciprocity of an unpolarized nearest-neighbor bond,
    \item parity evenness of the active vacuum.
\end{itemize}

Introduce the tangent contact coefficients
\begin{equation}
    \kappa_{\mathrm{n}}, \qquad \kappa_{\mathrm{t}}, \qquad \kappa_{\mathrm{b}}, \qquad \kappa_{\mathrm{tw}}, \qquad \kappa_\varphi.
\end{equation}

All coefficients are tangent coefficients of the active phase. Their dimensions are chosen so that the contact energy is expressed per unit participating interface area ($[\kappa] = \mathrm{N/m^3}$). The rotational and phase variables enter through the associated arc-length factors $r_{\mathrm{D}}\boldsymbol{\vartheta}$ and $r_{\mathrm{D}}\varphi$.

The parity-even quadratic contact energy is
\begin{equation}
\begin{aligned}
    \mathcal{U}_{ij}^{(2)} ={}& \frac{1}{2} \kappa_{\mathrm{n}} (\delta_{\mathrm{n}}^{(ij)})^2 + \frac{1}{2} \kappa_{\mathrm{t}} |\boldsymbol{\delta}_{\mathrm{t}}^{(ij)}|^2 + \frac{1}{2} \kappa_{\mathrm{b}} r_{\mathrm{D}}^2 |\boldsymbol{\delta}_{\mathrm{b}}^{(ij)}|^2 \\
    &+ \frac{1}{2} \kappa_{\mathrm{tw}} r_{\mathrm{D}}^2 (\delta_{\mathrm{tw}}^{(ij)})^2 + \frac{1}{2} \kappa_\varphi r_{\mathrm{D}}^2 (\Delta\varphi_{ij})^2.
\end{aligned}
\end{equation}

No tangential--bending cross term is included in the baseline active vacuum. The candidate scalar $\boldsymbol{\delta}_{\mathrm{t}}^{(ij)} \cdot \boldsymbol{\delta}_{\mathrm{b}}^{(ij)}$ is parity odd: $\boldsymbol{\delta}_{\mathrm{t}}$ is a polar vector, while $\boldsymbol{\delta}_{\mathrm{b}}$, being derived from microrotation, is axial. Their dot product is a pseudoscalar. It cannot appear in the parity-even constitutive energy of an unpolarized inversion-symmetric vacuum. This is why the baseline contact action strictly excludes a $\kappa_{\mathrm{tb}}$ coefficient.

A chiral extension may contain a term of the form $\chi_{ij} \kappa_{\mathrm{tb}} r_{\mathrm{D}} \boldsymbol{\delta}_{\mathrm{t}}^{(ij)} \cdot \boldsymbol{\delta}_{\mathrm{b}}^{(ij)}$, where $\chi_{ij}$ is an explicit pseudoscalar bond chirality. Such a term belongs to a polarized or chiral defect sector. It vanishes after averaging over an unpolarized active vacuum and is not part of the baseline FCC contact action.

The harmonic phase term may be understood as the tangent approximation to a periodic phase-contact potential,
\begin{equation}
    \mathcal{U}_{\varphi}^{(ij)} = \kappa_\varphi r_{\mathrm{D}}^2 \left[ 1-\cos(\Delta\varphi_{ij}) \right].
\end{equation}
For $|\Delta\varphi_{ij}| \ll 1$, this becomes
\begin{equation}
    \mathcal{U}_{\varphi}^{(ij)} = \frac{1}{2} \kappa_\varphi r_{\mathrm{D}}^2 (\Delta\varphi_{ij})^2 + \mathcal{O}(\Delta\varphi_{ij}^4).
\end{equation}

The local stability conditions of the uncoupled parity-even bond are
\begin{equation}
    \kappa_{\mathrm{n}}>0, \qquad \kappa_{\mathrm{t}}>0, \qquad \kappa_{\mathrm{b}}>0, \qquad \kappa_{\mathrm{tw}}>0, \qquad \kappa_\varphi>0.
\end{equation}

It is useful to define the contact-strain vector
\begin{equation}
    \boldsymbol{\xi}_{ij} =
    \begin{pmatrix}
        \delta_{\mathrm{n}}^{(ij)} \\
        \boldsymbol{\delta}_{\mathrm{t}}^{(ij)} \\
        r_{\mathrm{D}}\boldsymbol{\delta}_{\mathrm{b}}^{(ij)} \\
        r_{\mathrm{D}}\delta_{\mathrm{tw}}^{(ij)} \\
        r_{\mathrm{D}}\Delta\varphi_{ij}
    \end{pmatrix}.
\end{equation}

With the tangential and bending entries understood as two-component vectors in the contact tangent plane, the bond energy can be written compactly as $\mathcal{U}_{ij}^{(2)} = \frac{1}{2} \boldsymbol{\xi}_{ij}^{\mathsf{T}} \mathbf{H}_{\mathrm{c}} \boldsymbol{\xi}_{ij}$. The baseline parity-even contact tangent matrix is block diagonal:
\begin{equation}
    \mathbf{H}_{\mathrm{c}} =
    \begin{pmatrix}
        \kappa_{\mathrm{n}} & 0 & 0 & 0 & 0 \\
        0 & \kappa_{\mathrm{t}}\mathbf{I}_2 & 0 & 0 & 0 \\
        0 & 0 & \kappa_{\mathrm{b}}\mathbf{I}_2 & 0 & 0 \\
        0 & 0 & 0 & \kappa_{\mathrm{tw}} & 0 \\
        0 & 0 & 0 & 0 & \kappa_\varphi
    \end{pmatrix}.
\end{equation}

Let the generalized two-node coordinate be
\begin{equation}
    \mathbf{q}_{ij} =
    \begin{pmatrix}
        \mathbf{u}_i \\
        \boldsymbol{\vartheta}_i \\
        \varphi_i \\
        \mathbf{u}_j \\
        \boldsymbol{\vartheta}_j \\
        \varphi_j
    \end{pmatrix}.
\end{equation}

Every component of $\boldsymbol{\xi}_{ij}$ is linear in the generalized coordinates, so $\boldsymbol{\xi}_{ij} = \mathbf{B}_{ij}\mathbf{q}_{ij}$, where $\mathbf{B}_{ij}$ is the objective bond-kinematic operator. The quadratic bond energy becomes
\begin{equation}
    \mathcal{U}_{ij}^{(2)} = \frac{1}{2} \mathbf{q}_{ij}^{\mathsf{T}} \mathbf{K}_{ij} \mathbf{q}_{ij},
\end{equation}
with $\mathbf{K}_{ij} = \mathbf{B}_{ij}^{\mathsf{T}} \mathbf{H}_{\mathrm{c}} \mathbf{B}_{ij}$. This is the stiffness block contributed by one nearest-neighbor contact to the FCC assembly.

The bond is reciprocal under endpoint interchange ($i \leftrightarrow j$). The squared magnitudes in the energy are unchanged, and because the phase energy is even in $\Delta\varphi_{ij}$, the complete bond energy satisfies $\mathcal{U}_{ij}^{(2)} = \mathcal{U}_{ji}^{(2)}$.

For a mode with no normal deformation, no relative microrotation, and no phase difference, the contact energy reduces to
\begin{equation}
    \mathcal{U}_{ij}^{(2)} = \frac{1}{2} \kappa_{\mathrm{t}} |\boldsymbol{\delta}_{\mathrm{t}}^{(ij)}|^2.
\end{equation}
For the restricted directional transverse mode, consistency requires $\kappa_{\mathrm{t}} = \tau_{\mathrm{D}} / d_{\mathrm{nn}}^2$. Under the isotropic contact condition ($\kappa_{\mathrm{n}}=\kappa_{\mathrm{t}}$), the common contact coefficient becomes $\kappa_{\mathrm{c}} = \tau_{\mathrm{D}} / d_{\mathrm{nn}}^2$. This recovers the scalar transverse contact law as one polarization of the objective three-dimensional nearest-neighbor action.

\section{Primitive-Cell Kinetic Bookkeeping and Gravitationally Referred Inertia}

The bond stiffness matrix may be assembled over the six undirected nearest-neighbor contacts of a primitive FCC cell. The kinetic inertia, however, must not be assigned independently to each contact as a separate copy of the primitive node mass. Translational mass, microrotational inertia, phase inertia if present, and entrained condensate inertia belong to the primitive-cell coordinate system. The bond network contributes potential energy; the cell carries the shared kinetic energy.

The correct bookkeeping structure is
\begin{equation}
    U_{\mathrm{cell}} = \sum_{a=1}^{6}U_a,
\end{equation}
but
\begin{equation}
    T_{\mathrm{cell}}^{(0)} = \frac{1}{2} \dot{\mathbf{q}}_{\mathrm{cell}}^{\mathsf{T}} \mathbf{M}_{\mathrm{cell}}^{(0)} \dot{\mathbf{q}}_{\mathrm{cell}}.
\end{equation}

Here $a=1,\ldots,6$ labels the six independently counted undirected nearest-neighbor bonds of the FCC primitive cell. The primitive-cell generalized coordinate may be written schematically as
\begin{equation}
    \mathbf{q}_{\mathrm{cell}} =
    \begin{pmatrix}
        \mathbf{u} \\
        \boldsymbol{\vartheta} \\
        \varphi
    \end{pmatrix},
\end{equation}
with the understanding that the independent dynamical status of $\varphi$ is not fixed by the present appendix. 

The matrix $\mathbf{M}_{\mathrm{cell}}^{(0)}$ is a symmetric positive primitive-cell mass matrix. It includes the effective inertial participation of the coherent node--condensate transfer coordinate. It is not a sum of six independent full-node masses.

The endpoint-transfer constraint acts on the completed active-FCC coordinate relative to the latent node--condensate reconfiguration coordinate. At the primitive-cell level, represent this by a transfer operator
\begin{equation}
    \mathbf{q}_{\mathrm{cell}} = \mathbf{T}_{\mathrm{cell}} \mathbf{y}_{\mathrm{cell}}.
\end{equation}

The transfer operator acts on the translational transverse state-transfer sector through the endpoint factor \(\mathcal C\). For a Bloch mode with nonzero wave vector \(\mathbf k\), define the propagation direction
\begin{equation}
    \hat{\mathbf k}
    =
    \frac{\mathbf k}{|\mathbf k|}.
\end{equation}
The translational projector onto the subspace transverse to the propagation direction is
\begin{equation}
    \mathbf P_{\mathrm{tr}}(\hat{\mathbf k})
    =
    \mathbf I
    -
    \hat{\mathbf k}\hat{\mathbf k}^{\mathsf T}.
\end{equation}

This projector is distinct from the bond-level tangential projector
\begin{equation}
    \mathbf P_{\perp}^{(a)}
    =
    \mathbf I
    -
    \hat{\mathbf n}_{a}\hat{\mathbf n}_{a}^{\mathsf T},
\end{equation}
which projects onto the plane tangent to a particular FCC nearest-neighbor bond. The operator \(\mathbf P_{\mathrm{tr}}(\hat{\mathbf k})\) is transverse to the propagation direction, whereas \(\mathbf P_{\perp}^{(a)}\) is tangential to the \(a\)-th bond direction. These two projectors are not generally identical.

The primitive-cell generalized coordinate is
\begin{equation}
    \mathbf q_{\mathrm{cell}}
    =
    \begin{pmatrix}
        \mathbf u\\
        \boldsymbol\vartheta\\
        \varphi
    \end{pmatrix}.
\end{equation}
Let \(\mathbb P_{\mathrm{tr}}\) denote the embedding of \(\mathbf P_{\mathrm{tr}}(\hat{\mathbf k})\) into this full generalized coordinate, acting on the translational displacement block and vanishing on the microrotational and phase blocks:
\begin{equation}
    \mathbb P_{\mathrm{tr}}
    =
    \begin{pmatrix}
        \mathbf P_{\mathrm{tr}}(\hat{\mathbf k}) & 0 & 0\\
        0 & 0 & 0\\
        0 & 0 & 0
    \end{pmatrix}.
\end{equation}

The local transfer operator may then be represented as
\begin{equation}
    \mathbf T_{\mathrm{cell}}
    =
    \mathbf I
    +
    (\mathcal C-1)\mathbb P_{\mathrm{tr}}.
\end{equation}

Sectors not included in the endpoint-transfer reduction remain unscaled at this order unless a separate loading law is derived for them. Thus the normal translational component, microrotational components, and phase coordinate are not automatically assigned the same gravitational transfer factor. The only sector for which the present appendix derives the referred inertial law is the translational transverse completed-transfer sector.

Because the gravitational background is quasistatic ($\dot{\mathbf{T}}_{\mathrm{cell}}=0$), $\dot{\mathbf{y}}_{\mathrm{cell}} = \mathbf{T}_{\mathrm{cell}}^{-1} \dot{\mathbf{q}}_{\mathrm{cell}}$. The reduced kinetic energy in completed coordinates is
\begin{equation}
    T_{\mathrm{cell}} = \frac{1}{2} \dot{\mathbf{q}}_{\mathrm{cell}}^{\mathsf{T}} \mathbf{M}_{\mathrm{cell}}^{(\mathrm{eff})} \dot{\mathbf{q}}_{\mathrm{cell}},
\end{equation}
with
\begin{equation}
    \mathbf{M}_{\mathrm{cell}}^{(\mathrm{eff})} = \mathbf{T}_{\mathrm{cell}}^{-\mathsf{T}} \mathbf{M}_{\mathrm{cell}}^{(0)} \mathbf{T}_{\mathrm{cell}}^{-1}.
\end{equation}

On the transverse transfer sector,
\begin{equation}
    \mathbb{P}_{\mathrm{tr}} \mathbf{M}_{\mathrm{cell}}^{(\mathrm{eff})} \mathbb{P}_{\mathrm{tr}} = \frac{1}{\mathcal{C}^2} \mathbb{P}_{\mathrm{tr}} \mathbf{M}_{\mathrm{cell}}^{(0)} \mathbb{P}_{\mathrm{tr}}.
\end{equation}

For a locally uniform gravitational strain, $\mathcal{C}=(1-\alpha_{\mathrm{G}})^2$. Therefore,
\begin{equation}
    \mathbf{M}_{\mathrm{tr}}^{(\mathrm{eff})} = \frac{\mathbf{M}_{\mathrm{tr}}^{(0)}}{(1-\alpha_{\mathrm{G}})^4}.
\end{equation}

The potential energy is not multiplied by this factor. The active-phase bond Hessians remain
\begin{equation}
    \mathbf{K}_a(\alpha_{\mathrm{G}}) = \mathbf{K}_a(0), \qquad 0 \leq \alpha_{\mathrm{G}} < 1,
\end{equation}
at the retained constitutive order. Gravitational loading is therefore represented as a referred inertia of the completed transverse transfer coordinate, not as continuous softening of the local FCC contact stiffness.

The primitive-cell quadratic action is
\begin{equation}
    S_{\mathrm{cell}}^{(2)} = \int \left[ \frac{1}{2} \dot{\mathbf{q}}_{\mathrm{cell}}^{\mathsf{T}} \mathbf{M}_{\mathrm{cell}}^{(\mathrm{eff})} \dot{\mathbf{q}}_{\mathrm{cell}} - \frac{1}{2} \mathbf{q}_{\mathrm{cell}}^{\mathsf{T}} \mathbf{K}_{\mathrm{cell}} \mathbf{q}_{\mathrm{cell}} \right] \mathrm{d}t,
\end{equation}
where $\mathbf{K}_{\mathrm{cell}} = \sum_{a=1}^{6} \mathbf{K}_a$.

The corresponding linearized equation is
\begin{equation}
    \mathbf{M}_{\mathrm{cell}}^{(\mathrm{eff})} \ddot{\mathbf{q}}_{\mathrm{cell}} + \mathbf{K}_{\mathrm{cell}} \mathbf{q}_{\mathrm{cell}} = \mathbf{0}.
\end{equation}

This equation displays the correct assembly logic: the six nearest-neighbor bonds contribute stiffness, but the primitive cell carries the shared inertia once. The inverse-fourth gravitational referred-inertia law remains intact, but it is now correctly applied to the assembled transverse kinetic sector rather than independently to every bond.

The null-mode requirements are also preserved. For uniform translation, compatible rigid rotation, and uniform phase shift, the primitive-cell stiffness assembly retains the required zero-energy directions.

\section{Nearest-Neighbor FCC Assembly and Formal Bloch Stiffness}

The active Soft Floor is an FCC Bravais lattice with one generalized node per primitive cell. Each node has twelve nearest neighbors. Because opposite directions represent the same undirected contact, the primitive-cell potential energy is assembled from six independently counted nearest-neighbor bonds.

The nearest-neighbor separation is
\begin{equation}
    d_{\mathrm{nn}} = 2r_{\mathrm{D}}.
\end{equation}

A convenient set of six undirected nearest-neighbor vectors is
\begin{equation}
    \mathbf{d}_1 = \frac{d_{\mathrm{nn}}}{\sqrt{2}}(1,1,0), \qquad \mathbf{d}_2 = \frac{d_{\mathrm{nn}}}{\sqrt{2}}(1,-1,0),
\end{equation}
\begin{equation}
    \mathbf{d}_3 = \frac{d_{\mathrm{nn}}}{\sqrt{2}}(1,0,1), \qquad \mathbf{d}_4 = \frac{d_{\mathrm{nn}}}{\sqrt{2}}(1,0,-1),
\end{equation}
\begin{equation}
    \mathbf{d}_5 = \frac{d_{\mathrm{nn}}}{\sqrt{2}}(0,1,1), \qquad \mathbf{d}_6 = \frac{d_{\mathrm{nn}}}{\sqrt{2}}(0,1,-1).
\end{equation}

The corresponding unit vectors are $\hat{\mathbf{n}}_a = \mathbf{d}_a / d_{\mathrm{nn}}$ for $a=1,\ldots,6$. The negatives of these six directions supply the remaining six directed neighbors.

The conventional cubic FCC cell has edge length $a_{\mathrm{FCC}} = 2\sqrt{2}r_{\mathrm{D}}$. Since the conventional FCC cell contains four lattice nodes, the primitive-cell volume is
\begin{equation}
    V_p = \frac{a_{\mathrm{FCC}}^3}{4} = 4\sqrt{2}r_{\mathrm{D}}^3.
\end{equation}

Let the generalized primitive-cell coordinate at lattice position $\mathbf{R}$ be
\begin{equation}
    \mathbf{q}_{\mathbf{R}} =
    \begin{pmatrix}
        \mathbf{u}_{\mathbf{R}} \\
        \boldsymbol{\vartheta}_{\mathbf{R}} \\
        \varphi_{\mathbf{R}}
    \end{pmatrix}.
\end{equation}

For a Bloch mode, $\mathbf{q}_{\mathbf{R}}(t) = \mathbf{q}_0 \exp\left[i(\mathbf{k}\cdot\mathbf{R}-\omega t)\right]$. The neighboring cell connected through $\mathbf{d}_a$ has $\mathbf{q}_{\mathbf{R}+\mathbf{d}_a} = e^{i\mathbf{k}\cdot\mathbf{d}_a} \mathbf{q}_{\mathbf{R}}$.

For contact $a$, the objective bond deformation vector is linear in the Bloch amplitude: $\boldsymbol{\xi}_a = \mathbf{B}_a(\mathbf{k})\mathbf{q}_0$. Here $\mathbf{B}_a(\mathbf{k})$ is the Bloch contact-kinematic operator associated with the $a$-th FCC direction. It contains the phase factor $e^{i\mathbf{k}\cdot\mathbf{d}_a}$, the normal and transverse projectors, the microrotation terms, and the internal phase difference.

Using the parity-even contact tangent matrix $\mathbf{H}_c$, the stiffness contribution of direction $a$ is
\begin{equation}
    \mathbf{K}_a(\mathbf{k}) = \mathbf{B}_a^\dagger(\mathbf{k}) \mathbf{H}_c \mathbf{B}_a(\mathbf{k}).
\end{equation}

The primitive-cell Bloch stiffness is therefore formally
\begin{equation}
    \boxed{ \mathbf{K}_{\mathrm{FCC}}(\mathbf{k}) = \sum_{a=1}^{6} \mathbf{B}_a^\dagger(\mathbf{k}) \mathbf{H}_c \mathbf{B}_a(\mathbf{k}) }.
\end{equation}

This expression defines the nearest-neighbor FCC stiffness assembly. It is not yet the completed full dynamic matrix calculation. The explicit matrices $\mathbf{B}_a(\mathbf{k})$, the full finite-$\mathbf{k}$ spectrum, the rotational branches, the phase branches, and the stability proof over the complete first Brillouin zone remain open calculations. This appendix uses the assembly only to extract the long-wavelength acoustic tensor.

Let the unstrained primitive-cell mass matrix be $\mathbf{M}_{\mathrm{cell}}^{(0)}$. The formal unstrained Bloch eigenvalue problem is
\begin{equation}
    \det\left[ \mathbf{K}_{\mathrm{FCC}}(\mathbf{k}) - \omega^2\mathbf{M}_{\mathrm{cell}}^{(0)} \right] = 0.
\end{equation}

Under gravitational loading, the completed transverse sector of the mass matrix is modified by the endpoint-transfer operator, giving
\begin{equation}
    \mathbf{M}_{\mathrm{cell}}^{(\mathrm{eff})} = \mathbf{T}_{\mathrm{cell}}^{-\mathsf{T}} \mathbf{M}_{\mathrm{cell}}^{(0)} \mathbf{T}_{\mathrm{cell}}^{-1}.
\end{equation}

The corresponding loaded eigenvalue problem is
\begin{equation}
    \det\left[ \mathbf{K}_{\mathrm{FCC}}(\mathbf{k}) - \omega^2\mathbf{M}_{\mathrm{cell}}^{(\mathrm{eff})} \right] = 0.
\end{equation}

In the transverse transfer subspace,
\begin{equation}
    \mathbf{M}_{\mathrm{tr}}^{(\mathrm{eff})} = \frac{\mathbf{M}_{\mathrm{tr}}^{(0)}}{\mathcal{C}^2}.
\end{equation}

In the local continuum limit, $\mathcal{C}=(1-\alpha_{\mathrm{G}})^2$, so
\begin{equation}
    \mathbf{M}_{\mathrm{tr}}^{(\mathrm{eff})} = \frac{\mathbf{M}_{\mathrm{tr}}^{(0)}}{(1-\alpha_{\mathrm{G}})^4}.
\end{equation}

The stiffness assembly remains $\mathbf{K}_{\mathrm{FCC}}(\mathbf{k};\alpha_{\mathrm{G}}) = \mathbf{K}_{\mathrm{FCC}}(\mathbf{k};0)$ at the retained active-phase tangent order.

At zero wave number ($\mathbf{k}=\mathbf{0}$), uniform translation of the entire lattice produces no contact deformation. The stiffness matrix must therefore contain three translational null modes. A compatible rigid rotation also gives zero objective contact strain. A uniform internal phase shift gives $\Delta\varphi_{ij}=0$ for every bond and is also a zero-energy direction of the phase-difference contact energy. Whether this phase direction represents a gauge redundancy, a true gapless mode, a locally locked coordinate, or a defect-confined variable is not determined here.

Thus, the only spectrum statement made in this appendix is the following: the translational acoustic sector has two long-wavelength transverse branches governed by the shear tensor derived below. No claim is made here that the full active-vacuum spectrum contains only two branches.

\section{Long-Wavelength FCC Homogenization and Cubic Elastic Constants}

The leading continuum tensor is obtained from a spatially affine displacement field $\mathbf{u}(\mathbf{x}) = \mathbf{H}\mathbf{x}$, where $\mathbf{H}=\nabla\mathbf{u}$. Decompose $\mathbf{H}$ into symmetric and antisymmetric parts:
\begin{equation}
    \mathbf{H} = \boldsymbol{\varepsilon}+\mathbf{W},
\end{equation}
with $\boldsymbol{\varepsilon} = \frac{1}{2}(\mathbf{H}+\mathbf{H}^{\mathsf{T}})$ and $\mathbf{W} = \frac{1}{2}(\mathbf{H}-\mathbf{H}^{\mathsf{T}})$. The symmetric tensor $\boldsymbol{\varepsilon}$ is the infinitesimal strain. The antisymmetric tensor $\mathbf{W}$ represents local rigid rotation. Let $\boldsymbol{\omega}$ be the axial vector associated with $\mathbf{W}$, so that $\mathbf{W}\hat{\mathbf{n}} = \boldsymbol{\omega}\times\hat{\mathbf{n}}$.

Across contact $a$, the affine relative center displacement is $\mathbf{u}(\mathbf{x}+\mathbf{d}_a)-\mathbf{u}(\mathbf{x}) = d_{\mathrm{nn}} \mathbf{H}\hat{\mathbf{n}}_a$. For a spatially uniform microrotation $\boldsymbol{\vartheta}$, the objective relative contact-point displacement becomes
\begin{equation}
    \Delta\mathbf{u}_a^{\mathrm{c}} = d_{\mathrm{nn}} \left[ \boldsymbol{\varepsilon}\hat{\mathbf{n}}_a + (\boldsymbol{\omega}-\boldsymbol{\vartheta}) \times \hat{\mathbf{n}}_a \right].
\end{equation}

The normal deformation is therefore $\delta_{\mathrm{n}}^{(a)} = d_{\mathrm{nn}} \hat{\mathbf{n}}_a \cdot \boldsymbol{\varepsilon} \hat{\mathbf{n}}_a$. The tangential deformation is
\begin{equation}
    \boldsymbol{\delta}_{\mathrm{t}}^{(a)} = d_{\mathrm{nn}} \mathbf{P}_\perp^{(a)} \left[ \boldsymbol{\varepsilon}\hat{\mathbf{n}}_a + (\boldsymbol{\omega}-\boldsymbol{\vartheta}) \times \hat{\mathbf{n}}_a \right],
\end{equation}
where $\mathbf{P}_\perp^{(a)} = \mathbf{I}-\hat{\mathbf{n}}_a\hat{\mathbf{n}}_a^{\mathsf{T}}$.

The six FCC directions satisfy the second-order identity $\sum_{a=1}^{6} n_i^{(a)}n_j^{(a)} = 2\delta_{ij}$. For a symmetric strain tensor, this identity implies that the summed cross term between $\boldsymbol{\varepsilon}\hat{\mathbf{n}}_a$ and $(\boldsymbol{\omega}-\boldsymbol{\vartheta})\times\hat{\mathbf{n}}_a$ vanishes. The relative-rotation contribution is positive and proportional to $|\boldsymbol{\omega}-\boldsymbol{\vartheta}|^2$.

Minimizing the affine energy with respect to the uniform microrotation gives $\boldsymbol{\vartheta}=\boldsymbol{\omega}$. Thus the microrotation follows the macroscopic rigid rotation. After relaxation,
\begin{equation}
    \boldsymbol{\delta}_{\mathrm{t}}^{(a)} = d_{\mathrm{nn}} \mathbf{P}_\perp^{(a)} \boldsymbol{\varepsilon}\hat{\mathbf{n}}_a.
\end{equation}

For a homogeneous affine deformation, neighboring cells share the same microrotation and internal phase. Therefore, $\boldsymbol{\delta}_{\mathrm{b}}^{(a)}=\mathbf{0}$, $\delta_{\mathrm{tw}}^{(a)}=0$, and $\Delta\varphi_a=0$. The bending, twisting, and phase sectors do not contribute to the leading affine elastic tensor.

Let $A_{\mathrm{ct}}$ denote the effective modal area measure multiplying the contact energy per unit interface area. The leading scaffold strain-energy density is
\begin{equation}
    W_{\mathrm{FCC}} = \frac{A_{\mathrm{ct}}}{V_p} \sum_{a=1}^{6} \left[ \frac{1}{2}\kappa_{\mathrm{n}} (\delta_{\mathrm{n}}^{(a)})^2 + \frac{1}{2}\kappa_{\mathrm{t}} |\boldsymbol{\delta}_{\mathrm{t}}^{(a)}|^2 \right].
\end{equation}

Using the relaxed affine contact kinematics,
\begin{equation}
    |\boldsymbol{\delta}_{\mathrm{t}}^{(a)}|^2 = d_{\mathrm{nn}}^2 \left[ |\boldsymbol{\varepsilon}\hat{\mathbf{n}}_a|^2 - (\hat{\mathbf{n}}_a\cdot \boldsymbol{\varepsilon} \hat{\mathbf{n}}_a)^2 \right].
\end{equation}

Therefore,
\begin{equation}
    W_{\mathrm{FCC}} = \frac{A_{\mathrm{ct}}d_{\mathrm{nn}}^2}{2V_p} \sum_{a=1}^{6} \left[ \kappa_{\mathrm{t}} |\boldsymbol{\varepsilon}\hat{\mathbf{n}}_a|^2 + (\kappa_{\mathrm{n}}-\kappa_{\mathrm{t}}) (\hat{\mathbf{n}}_a\cdot \boldsymbol{\varepsilon} \hat{\mathbf{n}}_a)^2 \right].
\end{equation}

The required FCC directional sums are $\sum_{a=1}^{6} |\boldsymbol{\varepsilon}\hat{\mathbf{n}}_a|^2 = 2\boldsymbol{\varepsilon}:\boldsymbol{\varepsilon}$, and
\begin{equation}
\begin{aligned}
    \sum_{a=1}^{6} (\hat{\mathbf{n}}_a\cdot \boldsymbol{\varepsilon} \hat{\mathbf{n}}_a)^2 ={}& \varepsilon_{xx}^2 + \varepsilon_{yy}^2 + \varepsilon_{zz}^2 \\
    &+ \varepsilon_{xx}\varepsilon_{yy} + \varepsilon_{xx}\varepsilon_{zz} + \varepsilon_{yy}\varepsilon_{zz} \\
    &+ 2(\varepsilon_{xy}^2+\varepsilon_{xz}^2+\varepsilon_{yz}^2).
\end{aligned}
\end{equation}

Substitution gives
\begin{equation}
\begin{aligned}
    W_{\mathrm{FCC}} = \frac{A_{\mathrm{ct}}d_{\mathrm{nn}}^2}{2V_p} \Big[ & (\kappa_{\mathrm{n}}+\kappa_{\mathrm{t}}) (\varepsilon_{xx}^2+\varepsilon_{yy}^2+\varepsilon_{zz}^2) \\
    & + (\kappa_{\mathrm{n}}-\kappa_{\mathrm{t}}) (\varepsilon_{xx}\varepsilon_{yy} +\varepsilon_{xx}\varepsilon_{zz} +\varepsilon_{yy}\varepsilon_{zz}) \\
    & + 2(\kappa_{\mathrm{n}}+\kappa_{\mathrm{t}}) (\varepsilon_{xy}^2+\varepsilon_{xz}^2+\varepsilon_{yz}^2) \Big].
\end{aligned}
\end{equation}

The strain-energy density of a cubic continuum is
\begin{equation}
\begin{aligned}
    W_{\mathrm{cub}} ={}& \frac{1}{2} C_{11} (\varepsilon_{xx}^2+\varepsilon_{yy}^2+\varepsilon_{zz}^2) \\
    & + C_{12} (\varepsilon_{xx}\varepsilon_{yy} +\varepsilon_{xx}\varepsilon_{zz} +\varepsilon_{yy}\varepsilon_{zz}) \\
    & + 2C_{44} (\varepsilon_{xy}^2+\varepsilon_{xz}^2+\varepsilon_{yz}^2).
\end{aligned}
\end{equation}

Matching terms gives the leading FCC scaffold constants:
\begin{equation}
    C_{11}^{(\mathrm{FCC})} = \frac{A_{\mathrm{ct}}d_{\mathrm{nn}}^2}{V_p} (\kappa_{\mathrm{n}}+\kappa_{\mathrm{t}}),
\end{equation}
\begin{equation}
    C_{12}^{(\mathrm{FCC})} = \frac{A_{\mathrm{ct}}d_{\mathrm{nn}}^2}{2V_p} (\kappa_{\mathrm{n}}-\kappa_{\mathrm{t}}),
\end{equation}
\begin{equation}
    C_{44}^{(\mathrm{FCC})} = \frac{A_{\mathrm{ct}}d_{\mathrm{nn}}^2}{2V_p} (\kappa_{\mathrm{n}}+\kappa_{\mathrm{t}}).
\end{equation}

These constants define a cubic continuum. FCC geometry alone guarantees the cubic tensor structure, not exact isotropy. Exact isotropy requires a further relation among the tangent contact coefficients.

The leading cubic continuum is mechanically stable when $C_{44}>0$, $C_{11}-C_{12}>0$, and $C_{11}+2C_{12}>0$. For the FCC scaffold, this requires
\begin{equation}
    \kappa_{\mathrm{n}}>0, \qquad \kappa_{\mathrm{t}}>0.
\end{equation}
These conditions are sufficient for the stability of the leading affine scaffold continuum.

\section{Macroscopic Isotropy Condition of the Active FCC Scaffold}

For a cubic elastic continuum, the anisotropy factor is
\begin{equation}
    \mathcal{A}_{\mathrm{cub}} = \frac{2C_{44}}{C_{11}-C_{12}}.
\end{equation}

Using the FCC scaffold constants derived above,
\begin{equation}
    \mathcal{A}_{\mathrm{cub}} = \frac{ 2(\kappa_{\mathrm{n}}+\kappa_{\mathrm{t}}) }{ \kappa_{\mathrm{n}}+3\kappa_{\mathrm{t}} }.
\end{equation}

Exact isotropy requires $\mathcal{A}_{\mathrm{cub}}=1$. Therefore,
\begin{equation}
    2(\kappa_{\mathrm{n}}+\kappa_{\mathrm{t}}) = \kappa_{\mathrm{n}}+3\kappa_{\mathrm{t}},
\end{equation}
which yields
\begin{equation}
    \boxed{ \kappa_{\mathrm{n}} = \kappa_{\mathrm{t}} }.
\end{equation}

The active FCC scaffold is therefore isotropic at leading continuum order if and only if its incremental normal and tangential contact stiffnesses are equal.

Define the common contact coefficient $\kappa_{\mathrm{c}} \equiv \kappa_{\mathrm{n}} = \kappa_{\mathrm{t}}$. Then
\begin{equation}
    C_{11}^{(\mathrm{FCC})} = 2 \frac{A_{\mathrm{ct}}d_{\mathrm{nn}}^2}{V_p} \kappa_{\mathrm{c}}, \qquad C_{12}^{(\mathrm{FCC})}=0, \qquad C_{44}^{(\mathrm{FCC})} = \frac{A_{\mathrm{ct}}d_{\mathrm{nn}}^2}{V_p} \kappa_{\mathrm{c}}.
\end{equation}

Define the scaffold shear coefficient
\begin{equation}
    K_{\mathrm{sh}}^{(\mathrm{cell})} = \frac{A_{\mathrm{ct}}d_{\mathrm{nn}}^2}{V_p} \kappa_{\mathrm{c}}.
\end{equation}

Then the isotropic scaffold constants become
\begin{equation}
    C_{11}^{(\mathrm{FCC})} = 2K_{\mathrm{sh}}^{(\mathrm{cell})}, \qquad C_{12}^{(\mathrm{FCC})}=0, \qquad C_{44}^{(\mathrm{FCC})} = K_{\mathrm{sh}}^{(\mathrm{cell})}.
\end{equation}

The nearest-neighbor FCC scaffold therefore supplies isotropic shear rigidity but, by itself, does not supply a first Lamé coefficient. The isotropic tensor associated with this scaffold is
\begin{equation}
    C_{ijkl}^{(\mathrm{FCC})} = K_{\mathrm{sh}}^{(\mathrm{cell})} (\delta_{ik}\delta_{jl}+\delta_{il}\delta_{jk}).
\end{equation}

The volumetric part of the coherent active vacuum is not supplied by the scaffold alone. It belongs to the interstitial-condensate constraint and pressure sector. If a finite regularizing condensate coefficient $B_{\mathrm{IC}}$ is introduced temporarily, the complete isotropic tangent tensor takes the diagnostic form
\begin{equation}
    C_{ijkl} = B_{\mathrm{IC}}\delta_{ij}\delta_{kl} + K_{\mathrm{sh}}^{(\mathrm{cell})} (\delta_{ik}\delta_{jl}+\delta_{il}\delta_{jk}).
\end{equation}
However, in the active coherent branch, $B_{\mathrm{IC}}$ is not interpreted as an ordinary finite scalar acoustic modulus. It is a regularization of the incompressible constraint discussed separately in the following appendix.

The important result of the present appendix is the contact-level condition $\kappa_{\mathrm{n}}=\kappa_{\mathrm{t}}$. Cubic FCC geometry produces a cubic acoustic tensor. The equality of normal and tangential tangent contact stiffness is the additional constitutive condition that makes the leading continuum exactly isotropic.

\section{Effective Modal Normalization of the FCC Contact Coefficient}

The preceding homogenization expresses the leading scaffold shear coefficient as
\begin{equation}
    K_{\mathrm{sh}}^{(\mathrm{cell})} = \frac{A_{\mathrm{ct}}d_{\mathrm{nn}}^2}{V_p} \kappa_{\mathrm{c}}.
\end{equation}

The directional transverse contact law gives $\kappa_{\mathrm{c}} = \tau_{\mathrm{D}} / d_{\mathrm{nn}}^2$. Therefore,
\begin{equation}
    K_{\mathrm{sh}}^{(\mathrm{cell})} = \frac{A_{\mathrm{ct}}}{V_p} \tau_{\mathrm{D}}.
\end{equation}

The quantity $A_{\mathrm{ct}}$ must be interpreted carefully. It is not introduced here as a directly observed geometric patch on a Wigner–Seitz face. It is an effective modal area that converts the surface-inertia and surface-tension representation of the contact action into the primitive-cell translational normalization.

The Wigner–Seitz cell of an FCC lattice is a rhombic dodecahedron. Its twelve faces correspond to the twelve nearest-neighbor directions. For nearest-neighbor separation $d_{\mathrm{nn}}$, each rhombic face has diagonals $d_{\mathrm{nn}}$ and $d_{\mathrm{nn}} / \sqrt{2}$. The full Wigner–Seitz face area is therefore
\begin{equation}
    A_{\mathrm{WS}} = \frac{1}{2} d_{\mathrm{nn}} \frac{d_{\mathrm{nn}}}{\sqrt{2}} = \frac{d_{\mathrm{nn}}^2}{2\sqrt{2}}.
\end{equation}

Using $d_{\mathrm{nn}}=2r_{\mathrm{D}}$, one obtains $A_{\mathrm{WS}} = \sqrt{2}r_{\mathrm{D}}^2$. The ratio of one full Wigner–Seitz face to the primitive-cell volume is
\begin{equation}
    \frac{A_{\mathrm{WS}}}{V_p} = \frac{\sqrt{2}r_{\mathrm{D}}^2}{4\sqrt{2}r_{\mathrm{D}}^3} = \frac{1}{4r_{\mathrm{D}}}.
\end{equation}

This is a valid FCC geometric identity. It defines the full nearest-neighbor face area. It does not by itself identify the modal area associated with the reduced transverse transfer coordinate.

The modal area is fixed by matching the surface-inertia representation of the coherent transfer coordinate to the primitive-cell inertial normalization. Let $M_{\mathrm{D}}$ be the independently defined Soft Floor node inertia scale, and let $\mu_{\mathrm{D}}$ be the surface inertia. The modal area is defined by
\begin{equation}
    \boxed{ A_{\mathrm{ct}}^{(\mathrm{modal})} \equiv \frac{M_{\mathrm{D}}}{\mu_{\mathrm{D}}} }.
\end{equation}

This is a normalization identity, not a literal statement that each bond carries an independent mass $M_{\mathrm{D}}$. The primitive cell carries the shared translational inertia once. $A_{\mathrm{ct}}^{(\mathrm{modal})}$ is the area factor that makes the surface kinetic coefficient reproduce that primitive-cell inertia when expressed in the reduced coherent transfer coordinate.

Using the established TFT relations $M_{\mathrm{D}} = 2\pi\hbar / (c_{\mathrm{E}}r_{\mathrm{D}})$ and $\mu_{\mathrm{D}} = c_{\mathrm{E}}^2 / (G r_{\mathrm{D}})$, one obtains
\begin{equation}
    A_{\mathrm{ct}}^{(\mathrm{modal})} = \frac{ 2\pi\hbar/(c_{\mathrm{E}}r_{\mathrm{D}}) }{ c_{\mathrm{E}}^2/(G r_{\mathrm{D}}) } = \frac{2\pi\hbar G}{c_{\mathrm{E}}^3}.
\end{equation}

Since the substrate scale satisfies $\ell_{\mathrm{S}}^2 = \hbar G / c_{\mathrm{E}}^3$, the modal area is
\begin{equation}
    \boxed{ A_{\mathrm{ct}}^{(\mathrm{modal})} = 2\pi\ell_{\mathrm{S}}^2 }.
\end{equation}

The corresponding coherent transverse modal projection is defined by
\begin{equation}
    \iota_{\mathrm{tr}}^{(\mathrm{modal})}
    \equiv
    \frac{
    A_{\mathrm{ct}}^{(\mathrm{modal})}
    }{
    V_p
    }.
\end{equation}
Using
\begin{equation}
    V_p
    =
    4\sqrt2\,r_{\mathrm D}^3,
\end{equation}
and
\begin{equation}
    A_{\mathrm{ct}}^{(\mathrm{modal})}
    =
    2\pi\ell_{\mathrm S}^2,
\end{equation}
one obtains
\begin{equation}
    \iota_{\mathrm{tr}}^{(\mathrm{modal})}
    =
    \frac{
    \pi\ell_{\mathrm S}^2
    }{
    2\sqrt2\,r_{\mathrm D}^3
    }.
\end{equation}

This projection also satisfies
\begin{equation}
    \iota_{\mathrm{tr}}^{(\mathrm{modal})}
    =
    \frac{
    \rho_{\mathrm{eff}}^{(\mathrm{micro})}
    }{
    \mu_{\mathrm D}
    }.
\end{equation}
Indeed, since
\begin{equation}
    \rho_{\mathrm{eff}}^{(\mathrm{micro})}
    =
    \frac{M_{\mathrm D}}{V_p},
\end{equation}
and
\begin{equation}
    A_{\mathrm{ct}}^{(\mathrm{modal})}
    =
    \frac{M_{\mathrm D}}{\mu_{\mathrm D}},
\end{equation}
it follows that
\begin{equation}
    \frac{
    A_{\mathrm{ct}}^{(\mathrm{modal})}
    }{
    V_p
    }
    =
    \frac{
    \rho_{\mathrm{eff}}^{(\mathrm{micro})}
    }{
    \mu_{\mathrm D}
    }.
\end{equation}

The quantity \(\iota_{\mathrm{tr}}^{(\mathrm{modal})}\) is the coherent transverse projection associated with the periodic FCC transfer mode. It is not the localized defect-flow projection \(\iota_V^{(\mathrm{def})}\) introduced in Appendix~\ref{app:condensate_impedance}. The former belongs to the entrained periodic transverse acoustic sector; the latter belongs to a boundary-value problem for relative condensate flow around a localized defect.

The equivalent circular radius associated with this modal area is $r_{\mathrm{ct}}^{(\mathrm{modal})} = \sqrt{A_{\mathrm{ct}}^{(\mathrm{modal})} / \pi} = \sqrt{2}\ell_{\mathrm{S}}$. This result should be read as follows: $A_{\mathrm{ct}}^{(\mathrm{modal})}$ is the surface-to-cell normalization area of the coherent transfer mode. It should not be read as having been geometrically located as a literal physical patch on every FCC face.

The ratio to the full Wigner–Seitz face may still be useful as a dimensionless comparison:
\begin{equation}
    f_{\mathrm{ct}}^{(\mathrm{modal})} = \frac{A_{\mathrm{ct}}^{(\mathrm{modal})}}{A_{\mathrm{WS}}} = \sqrt{2}\pi \left(\frac{\ell_{\mathrm{S}}}{r_{\mathrm{D}}}\right)^2.
\end{equation}

Because $\ell_{\mathrm{S}}\ll r_{\mathrm{D}}$, $f_{\mathrm{ct}}^{(\mathrm{modal})}\ll1$. This ratio is an effective modal-area ratio. It does not prove that only this fraction of the Wigner–Seitz face is physically active as a literal geometric contact footprint. The full Wigner–Seitz face remains the geometric nearest-neighbor boundary; the modal area is the normalization measure of the reduced transverse transfer coordinate.

The integrated modal stiffness associated with the reduced transfer coordinate is $k_{\mathrm{ct}}^{(\mathrm{modal})} = A_{\mathrm{ct}}^{(\mathrm{modal})} \kappa_{\mathrm{c}}$. Using $\kappa_{\mathrm{c}} = \tau_{\mathrm{D}} / d_{\mathrm{nn}}^2$, one obtains
\begin{equation}
    k_{\mathrm{ct}}^{(\mathrm{modal})} = \frac{\pi}{2} \tau_{\mathrm{D}} \left(\frac{\ell_{\mathrm{S}}}{r_{\mathrm{D}}}\right)^2.
\end{equation}

The discrete directional speed associated with the reduced contact sequence is $c_T^{(0)} = d_{\mathrm{nn}} \sqrt{ k_{\mathrm{ct}}^{(\mathrm{modal})} / m_{\mathrm{ct}}^{(\mathrm{modal})} }$. Here $m_{\mathrm{ct}}^{(\mathrm{modal})} = A_{\mathrm{ct}}^{(\mathrm{modal})}\mu_{\mathrm{D}} = M_{\mathrm{D}}$. This equality is a modal normalization of the representative transfer coordinate, not a statement that each of the six primitive-cell bonds contains a separate full node mass.

Using $k_{\mathrm{ct}}^{(\mathrm{modal})}$, $m_{\mathrm{ct}}^{(\mathrm{modal})}$, and $\kappa_{\mathrm{c}}$, one finds
\begin{equation}
    (c_T^{(0)})^2 = d_{\mathrm{nn}}^2 \frac{ A_{\mathrm{ct}}^{(\mathrm{modal})}\kappa_{\mathrm{c}} }{ A_{\mathrm{ct}}^{(\mathrm{modal})}\mu_{\mathrm{D}} } = \frac{\tau_{\mathrm{D}}}{\mu_{\mathrm{D}}} = c_{\mathrm{E}}^2.
\end{equation}
Therefore, $c_T^{(0)}=c_{\mathrm{E}}$. The modal area cancels from the unstrained wave-speed normalization. Its role is not to fit $c_{\mathrm{E}}$, but to keep the surface and primitive-cell representations of the coherent transfer coordinate mutually consistent.

The primitive-cell microscopic inertia density associated with this bare cell-level tangent normalization is $\rho_{\mathrm{eff}}^{(\mathrm{micro})} = M_{\mathrm{D}} / V_p$. Using $M_{\mathrm{D}} = \rho_{\mathrm{D}} \left(\frac{4}{3}\pi r_{\mathrm{D}}^3\right)$ and $V_p=4\sqrt{2}r_{\mathrm{D}}^3$, one obtains
\begin{equation}
    \rho_{\mathrm{eff}}^{(\mathrm{micro})} = \frac{\pi}{3\sqrt{2}}\rho_{\mathrm{D}}.
\end{equation}

The FCC close-packing fraction is $\beta_{\mathrm{geo}} = \pi / (3\sqrt{2})$. Thus,
\begin{equation}
    \boxed{ \rho_{\mathrm{eff}}^{(\mathrm{micro})} = \beta_{\mathrm{geo}}\rho_{\mathrm{D}} }.
\end{equation}

The corresponding primitive-cell microscopic tangent shear coefficient is
\begin{equation}
    \boxed{ K_{\mathrm{sh}}^{(\mathrm{micro})} = \rho_{\mathrm{eff}}^{(\mathrm{micro})}c_{\mathrm{E}}^2 = \beta_{\mathrm{geo}}\rho_{\mathrm{D}}c_{\mathrm{E}}^2 }.
\end{equation}
Equivalently, $K_{\mathrm{sh}}^{(\mathrm{micro})} = (A_{\mathrm{ct}}^{(\mathrm{modal})} / V_p) \tau_{\mathrm{D}}$.

This coefficient is the bare primitive-cell tangent coefficient associated with the reduced coherent transfer coordinate. The associated inverse-length projection is \(\iota_{\mathrm{tr}}^{(\mathrm{modal})}\), not the localized defect projection \(\iota_V^{(\mathrm{def})}\). It must therefore not be automatically identified either with the defect-flow impedance calibration or with the macroscopic effective shear modulus inferred from long-baseline gravitational-wave propagation. The cross-scale distinction is addressed next.

\section{Microscopic Cell Coefficients versus Macroscopic Effective Moduli}

The primitive-cell construction fixes the bare microscopic tangent normalization of the active FCC transfer coordinate. In that normalization, the effective cell inertia density is
\begin{equation}
    \rho_{\mathrm{eff}}^{(\mathrm{micro})} = \beta_{\mathrm{geo}}\rho_{\mathrm{D}},
\end{equation}
and the corresponding microscopic shear coefficient is
\begin{equation}
    K_{\mathrm{sh}}^{(\mathrm{micro})} = \beta_{\mathrm{geo}}\rho_{\mathrm{D}}c_{\mathrm{E}}^2.
\end{equation}

These quantities are local tangent coefficients of the primitive active-FCC cell. They characterize the coherent transfer coordinate before macroscopic averaging, long-baseline projection, or gravitational-wave-scale matching has been applied. They should therefore not be identified directly with the large-scale effective moduli inferred from gravitational-wave propagation.

The manuscript also employs macroscopic gravitational-wave effective quantities of the approximate form
\begin{equation}
    K_{\mathrm{shear}}^{(\mathrm{macro})} \sim 9.4\times10^{30}\ \mathrm{Pa},
\end{equation}
and
\begin{equation}
    \rho_{\mathrm{eff}}^{(\mathrm{macro})} \sim 1.04\times10^{14}\ \mathrm{kg\,m^{-3}}.
\end{equation}

The primitive-cell normalization instead gives quantities of order
\begin{equation}
    K_{\mathrm{sh}}^{(\mathrm{micro})} \sim 10^{79}\ \mathrm{Pa},
\end{equation}
and
\begin{equation}
    \rho_{\mathrm{eff}}^{(\mathrm{micro})} \sim 10^{62}\ \mathrm{kg\,m^{-3}}.
\end{equation}

The difference is not a dimensional contradiction. Both the microscopic and macroscopic pairs preserve the same propagation-ratio condition,
\begin{equation}
    \frac{K}{\rho} = c_{\mathrm{E}}^2.
\end{equation}

The discrepancy lies in the absolute normalization of the participating stiffness and inertia. The microscopic coefficients describe the bare primitive-cell tangent response, whereas the macroscopic coefficients describe the effective continuum response sampled by coherent gravitational-wave propagation.

To keep these levels distinct, define the diagnostic microscopic-to-macroscopic matching ratio
\begin{equation}
    \boxed{ \zeta_{\mathrm{GW}} \equiv \frac{ K_{\mathrm{shear}}^{(\mathrm{macro})} }{ K_{\mathrm{sh}}^{(\mathrm{micro})} } = \frac{ \rho_{\mathrm{eff}}^{(\mathrm{macro})} }{ \rho_{\mathrm{eff}}^{(\mathrm{micro})} } }.
\end{equation}

With the representative manuscript values,
\begin{equation}
    \zeta_{\mathrm{GW}} \sim 10^{-48}.
\end{equation}

The macroscopic effective coefficients are then written as
\begin{equation}
    K_{\mathrm{shear}}^{(\mathrm{macro})} = \zeta_{\mathrm{GW}} K_{\mathrm{sh}}^{(\mathrm{micro})},
\end{equation}
and
\begin{equation}
    \rho_{\mathrm{eff}}^{(\mathrm{macro})} = \zeta_{\mathrm{GW}} \rho_{\mathrm{eff}}^{(\mathrm{micro})}.
\end{equation}

The equality of the stiffness and inertia ratios is a matching condition required to preserve the propagation normalization,
\begin{equation}
    \frac{ K_{\mathrm{shear}}^{(\mathrm{macro})} }{ \rho_{\mathrm{eff}}^{(\mathrm{macro})} } = \frac{ K_{\mathrm{sh}}^{(\mathrm{micro})} }{ \rho_{\mathrm{eff}}^{(\mathrm{micro})} } = c_{\mathrm{E}}^2.
\end{equation}

Thus $\zeta_{\mathrm{GW}}$ does not alter the wave-speed law. It records the distinction between the microscopic tangent normalization and the macroscopic effective continuum normalization while preserving their common stiffness-to-inertia ratio.

The ratio $\zeta_{\mathrm{GW}}$ is not derived in this appendix. It is introduced as a diagnostic matching ratio between two levels of description. A complete theory must derive it from the full multiscale FCC elastodynamic routing problem, including the coupling between primitive-cell contact mechanics and the long-baseline coherent shear mode sampled by astrophysical gravitational waves.

Until that derivation is supplied, the terminology must remain explicit:
\begin{equation}
    K_{\mathrm{sh}}^{(\mathrm{micro})} = \beta_{\mathrm{geo}}\rho_{\mathrm{D}}c_{\mathrm{E}}^2
\end{equation}
is the primitive-cell microscopic tangent shear coefficient, whereas $K_{\mathrm{shear}}^{(\mathrm{macro})}$ is the large-scale effective shear modulus used in the gravitational-wave continuum.

Similarly,
\begin{equation}
    \rho_{\mathrm{eff}}^{(\mathrm{micro})} = \beta_{\mathrm{geo}}\rho_{\mathrm{D}}
\end{equation}
is the primitive-cell microscopic inertial density, whereas $\rho_{\mathrm{eff}}^{(\mathrm{macro})}$ is the large-scale effective inertial density of the propagating shear disturbance.

This distinction also fixes the interpretation of gravitationally referred inertia. At any chosen effective scale, the endpoint-transfer law modifies the completed-coordinate inertia as
\begin{equation}
    \rho_{\mathrm{eff}}(\alpha_{\mathrm{G}}) = \frac{ \rho_{\mathrm{eff}}^{(0)} }{ (1-\alpha_{\mathrm{G}})^4 }.
\end{equation}

For a primitive-cell calculation,
\begin{equation}
    \rho_{\mathrm{eff}}^{(0)} = \rho_{\mathrm{eff}}^{(\mathrm{micro})}.
\end{equation}

For a gravitational-wave-scale continuum calculation,
\begin{equation}
    \rho_{\mathrm{eff}}^{(0)} = \rho_{\mathrm{eff}}^{(\mathrm{macro})}.
\end{equation}

Because the corresponding shear coefficient is matched at the same descriptive scale, the ratio $K/\rho$ remains invariant and the transverse propagation law remains
\begin{equation}
    c_P = c_{\mathrm{E}} (1-\alpha_{\mathrm{G}})^2.
\end{equation}

The microscopic and macroscopic descriptions therefore differ in absolute effective normalization, not in the local endpoint-transfer mechanism.

The primitive-cell values must not be promoted to unique macroscopic vacuum moduli. Conversely, the macroscopic gravitational-wave values must not be read back as primitive-cell contact coefficients. The two descriptions are connected by the diagnostic matching ratio $\zeta_{\mathrm{GW}}$, whose physical derivation remains an open multiscale problem.

\section{Recovery of the Constrained Transverse Propagation Law}

The preceding sections provide the ingredients required to recover the gravitationally constrained transverse propagation law in the full long-wavelength FCC setting.

The endpoint-transfer closure is
\begin{equation}
    \mathcal{C}_{ij} = (1-\alpha_{\mathrm{G}}^{(i)}) (1-\alpha_{\mathrm{G}}^{(j)}).
\end{equation}

For a gravitational loading varying negligibly over one nearest-neighbor separation, this reduces to $\mathcal{C} = (1-\alpha_{\mathrm{G}})^2$.

The constrained kinetic reduction gives the effective transverse inertia
\begin{equation}
    \mathbf{M}_{\mathrm{tr}}^{(\mathrm{eff})} = \frac{\mathbf{M}_{\mathrm{tr}}^{(0)}}{\mathcal{C}^2}.
\end{equation}

Equivalently, in density form,
\begin{equation}
    \rho_{\mathrm{eff}}(\alpha_{\mathrm{G}}) = \frac{\rho_{\mathrm{eff}}^{(0)}}{\mathcal{C}^2}.
\end{equation}

The active-phase contact stiffness is unchanged at the retained tangent order,
\begin{equation}
    K_{\mathrm{sh}}(\alpha_{\mathrm{G}}) = K_{\mathrm{sh}}(0), \qquad 0 \leq \alpha_{\mathrm{G}} < 1.
\end{equation}

The long-wavelength transverse acoustic equation therefore has the form
\begin{equation}
    \omega_T^2 = \frac{K_{\mathrm{sh}}(0)}{\rho_{\mathrm{eff}}(\alpha_{\mathrm{G}})} k^2 + \mathcal{O}(k^4d_{\mathrm{nn}}^2).
\end{equation}

Substituting the referred inertia gives
\begin{equation}
    \omega_T^2 = \frac{K_{\mathrm{sh}}(0)}{\rho_{\mathrm{eff}}^{(0)}} \mathcal{C}^2k^2 + \mathcal{O}(k^4d_{\mathrm{nn}}^2).
\end{equation}

The unstrained normalization is
\begin{equation}
    \frac{K_{\mathrm{sh}}(0)}{\rho_{\mathrm{eff}}^{(0)}} = c_{\mathrm{E}}^2.
\end{equation}

Hence,
\begin{equation}
    \omega_T^2 = c_{\mathrm{E}}^2 \mathcal{C}^2k^2 + \mathcal{O}(k^4d_{\mathrm{nn}}^2).
\end{equation}

Using the local continuum endpoint factor, $\mathcal{C}=(1-\alpha_{\mathrm{G}})^2$, one obtains
\begin{equation}
    \boxed{ \omega_T^2 = c_{\mathrm{E}}^2 (1-\alpha_{\mathrm{G}})^4 k^2 + \mathcal{O}(k^4d_{\mathrm{nn}}^2) }.
\end{equation}

Therefore the long-wavelength transverse propagation speed is
\begin{equation}
    \boxed{ c_P = c_{\mathrm{E}}(1-\alpha_{\mathrm{G}})^2 }.
\end{equation}

This result is the FCC-homogenized version of the directional contact-chain law. It follows through the sequence
\begin{equation}
    \text{endpoint capacity} \longrightarrow \mathcal{C} \longrightarrow \text{referred inertia} \longrightarrow \text{long-wavelength FCC tensor} \longrightarrow c_P.
\end{equation}

The result applies to the translational transverse acoustic sector of the active FCC phase. It does not classify the independent dynamical status of microrotational, torsional, internal phase, or localized defect modes.

In the active-side saturation limit, $\alpha_{\mathrm{G}} \rightarrow 1^-$, the completed-transfer speed tends to zero ($c_P \rightarrow 0$). The mechanism is not a vanishing of the baseline FCC contact stiffness. At the retained constitutive order, $K_{\mathrm{sh}}(\alpha_{\mathrm{G}}) = K_{\mathrm{sh}}(0)$, while the completed-coordinate inertia diverges:
\begin{equation}
    \rho_{\mathrm{eff}}(\alpha_{\mathrm{G}}) = \frac{\rho_{\mathrm{eff}}^{(0)}}{(1-\alpha_{\mathrm{G}})^4} \rightarrow \infty.
\end{equation}
Thus the active transverse mode is arrested by divergent referred inertia. The divergence indicates the loss of validity of the active FCC coordinate at the phase boundary; it is not interpreted as an infinite ordinary material density of the SC substrate.

For a nonuniform gravitational loading across a bond, the endpoint-resolved form is
\begin{equation}
    c_P^{(ij)} = c_{\mathrm{E}} (1-\alpha_{\mathrm{G}}^{(i)}) (1-\alpha_{\mathrm{G}}^{(j)}).
\end{equation}

Let $\alpha_{\mathrm{G}} = (\alpha_{\mathrm{G}}^{(i)}+\alpha_{\mathrm{G}}^{(j)}) / 2$ and $\delta\alpha_{\mathrm{G}} = \alpha_{\mathrm{G}}^{(j)}-\alpha_{\mathrm{G}}^{(i)}$. Then
\begin{equation}
    (1-\alpha_{\mathrm{G}}^{(i)}) (1-\alpha_{\mathrm{G}}^{(j)}) = (1-\alpha_{\mathrm{G}})^2 - \frac{(\delta\alpha_{\mathrm{G}})^2}{4}.
\end{equation}

The first-order endpoint-gradient correction cancels. The local continuum expression therefore neglects only a second-order endpoint-difference term. This gives the continuum propagation law a controlled long-wavelength error structure.

\section{Status of the Full FCC Dynamic Spectrum and Remaining Closures}

This appendix establishes the leading reduced contact mechanics and long-wavelength transverse acoustic tensor of the active FCC Soft Floor. The results are sufficient to connect the endpoint-transfer law to a conservative contact action and to the continuum transverse propagation speed. They do not constitute a complete finite-wave-number microscopic elastodynamic theory.

The first completed result is the objective nearest-neighbor contact kinematics. The local contact deformation is determined by the contact-point displacement
\begin{equation}
    \Delta\mathbf{u}_{ij}^{\mathrm{c}} = \mathbf{u}_j-\mathbf{u}_i - r_{\mathrm{D}} (\boldsymbol{\vartheta}_i+\boldsymbol{\vartheta}_j) \times \hat{\mathbf{n}}_{ij},
\end{equation}
rather than by the relative center displacement alone. This construction yields invariant contact measures that vanish under rigid translation, compatible rigid rotation, and uniform internal phase displacement.

The second completed result is the parity-even quadratic contact action. The baseline active vacuum admits normal, tangential, bending, twisting, and phase-difference quadratic terms. It excludes the tangential–bending pseudoscalar unless an explicit chiral pseudoscalar is introduced. The unpolarized active vacuum therefore contains no baseline term proportional to $\boldsymbol{\delta}_{\mathrm{t}}^{(ij)} \cdot \boldsymbol{\delta}_{\mathrm{b}}^{(ij)}$.

The third completed result is the work-conjugate reduced construction of the endpoint-transfer factor,
\begin{equation}
    \mathcal{C}_{ij} = (1-\alpha_{\mathrm{G}}^{(i)}) (1-\alpha_{\mathrm{G}}^{(j)}).
\end{equation}
This relation is used as a reduced constitutive closure. The appendix shows that it is compatible with a conservative constrained kinetic action, but does not claim that the endpoint product has been derived from a complete parent contact potential.

The fourth completed result is the gravitationally referred transverse inertia,
\begin{equation}
    \mathbf{M}_{\mathrm{tr}}^{(\mathrm{eff})} = \frac{\mathbf{M}_{\mathrm{tr}}^{(0)}}{\mathcal{C}^2}.
\end{equation}

The fifth completed result is the affine FCC homogenization. The six nearest-neighbor directions produce a cubic continuum. Exact long-wavelength isotropy follows only when $\kappa_{\mathrm{n}} = \kappa_{\mathrm{t}}$.

The sixth completed result is the effective modal normalization,
\begin{equation}
    A_{\mathrm{ct}}^{(\mathrm{modal})} = \frac{M_{\mathrm{D}}}{\mu_{\mathrm{D}}} = 2\pi\ell_{\mathrm{S}}^2.
\end{equation}
This quantity is a surface-to-cell normalization area for the coherent transfer coordinate. It is not interpreted as a directly established literal contact patch on each Wigner–Seitz face.

The seventh established result is the distinction between microscopic and macroscopic effective coefficients. The primitive-cell tangent normalization gives
\begin{equation}
    \rho_{\mathrm{eff}}^{(\mathrm{micro})} = \beta_{\mathrm{geo}}\rho_{\mathrm{D}}, \qquad K_{\mathrm{sh}}^{(\mathrm{micro})} = \beta_{\mathrm{geo}}\rho_{\mathrm{D}}c_{\mathrm{E}}^2.
\end{equation}
These coefficients must be distinguished from macroscopic gravitational-wave effective quantities. The bridge between the primitive-cell tangent normalization and the macroscopic gravitational-wave continuum is represented by the diagnostic microscopic-to-macroscopic matching ratio
\begin{equation}
    \zeta_{\mathrm{GW}} \equiv \frac{ K_{\mathrm{shear}}^{(\mathrm{macro})} }{ K_{\mathrm{sh}}^{(\mathrm{micro})} } = \frac{ \rho_{\mathrm{eff}}^{(\mathrm{macro})} }{ \rho_{\mathrm{eff}}^{(\mathrm{micro})} }.
\end{equation}
The equality of the stiffness and inertia ratios is imposed to preserve
\begin{equation}
    \frac{K}{\rho} = c_{\mathrm{E}}^2
\end{equation}
between the microscopic and macroscopic effective descriptions. The derivation of $\zeta_{\mathrm{GW}}$ from the full microscopic FCC routing problem remains open.

The eighth established result is the transverse acoustic law,
\begin{equation}
    \omega_T^2 = c_{\mathrm{E}}^2 (1-\alpha_{\mathrm{G}})^4 k^2 + \mathcal{O}(k^4 d_{\mathrm{nn}}^2),
\end{equation}
with
\begin{equation}
    c_P = c_{\mathrm{E}}(1-\alpha_{\mathrm{G}})^2.
\end{equation}
These results define the leading reduced microphysical foundation of coherent transverse propagation in the active FCC Soft Floor.

Several microscopic completions remain outside the scope of this appendix. A complete parent potential must still be constructed to derive the single-endpoint map
\begin{equation}
    u_i = \mathcal{C}_i y_i, \qquad \mathcal{C}_i = 1-\alpha_{\mathrm{G}}^{(i)},
\end{equation}
from normal contact deformation, transverse shear, phase mismatch, condensate response, and vertical Tension-Column loading. The present appendix supplies the reduced endpoint closure that such a parent potential must reproduce.

The full Bloch dynamic matrix also remains to be computed. The formal stiffness assembly is given, but the explicit matrices $\mathbf{B}_a(\mathbf{k})$, the full finite-$\mathbf{k}$ eigenvalue spectrum, and the rotational, torsional, and phase-sector branches are not evaluated here.

Likewise, the affine stability conditions do not constitute a global lattice-stability proof. The leading continuum stability follows from the stronger bond-level positivity conditions on the tangent contact coefficients, but the complete theory must still show
\begin{equation}
    \omega_s^2(\mathbf{k}) \geq 0
\end{equation}
for every branch $s$ and every wave vector $\mathbf{k}$ in the first Brillouin zone.

The internal phase sector also remains to be classified. Since the contact energy depends on $\Delta\varphi_{ij}$, a uniform phase displacement is a zero-energy direction of the bond energy. Whether this direction is gauge-redundant, locally locked, gapped, defect-confined, or dynamically active through an independent phase inertia is not determined by the present transverse homogenization.

The chiral extension is similarly deferred. Parity-odd terms may enter only with an explicit pseudoscalar chirality and must vanish in the unpolarized vacuum average.

Finally, the scale bridge between the primitive-cell tangent coefficients and the macroscopic gravitational-wave effective coefficients remains an open multiscale problem. The diagnostic matching ratio $\zeta_{\mathrm{GW}}$ must ultimately be derived from the full microscopic routing structure of the active FCC lattice.

The interstitial-condensate pressure constraint, the absence of an ordinary finite-frequency longitudinal acoustic scalar branch in the coherent sector, and the localized defect-induced relative-flow impedance are developed separately in Appendix \ref{app:condensate_impedance}. Those mechanisms are not included in the present FCC contact homogenization.

The status of the present appendix is therefore precise: it establishes the leading reduced FCC contact action, the objective bond kinematics, the affine acoustic tensor, the isotropy condition, and the constrained transverse propagation law. It does not complete the parent microscopic potential, the full Bloch spectrum, the phase-sector classification, or the multiscale projection to macroscopic gravitational-wave moduli.

